\input{preamble.tex}

\title{Отчет по лабораторной работе № 6. \\ Статическая маршрутизация VLAN}
\author{Сущенко Алина Николаевна \\ НПИбд-01-23}
\date{2026}

\begin{document}

\maketitle
\newpage

\tableofcontents

\newpage
\section{Цель работы}
Настроить статическую маршрутизацию VLAN в сети.

\newpage
\section{Выполнение работы}

\subsection{Топология сети}
В логической области проекта разместили маршрутизатор Cisco 2811 (msk-ansusenko-gw-1), подключили его к порту 24 коммутатора msk-ansusenko-sw-1 в соответствии с таблицей портов. (Рис. \ref{01}):

\begin{center}
    \centering
    \includegraphics[width=\textwidth]{../images/01.png}
    \captionof{figure}{Обновлённая топология сети.}
    \label{01}
\end{center}

\subsection{Первоначальная настройка маршрутизатора}

На маршрутизаторе выполнили первоначальную настройку: задали имя, пароли, настроили удаленное подключение по SSH (Рис. \ref{02}):

\begin{minted}{cisco}
Router>enable
Router#configure terminal
Router(config)#hostname msk-ansusenko-gw-1
msk-ansusenko-gw-1(config)#line vty 0 4
msk-ansusenko-gw-1(config-line)#password cisco
msk-ansusenko-gw-1(config-line)#login
msk-ansusenko-gw-1(config-line)#line console 0
msk-ansusenko-gw-1(config-line)#password cisco
msk-ansusenko-gw-1(config-line)#login
msk-ansusenko-gw-1(config)#enable secret cisco
msk-ansusenko-gw-1(config)#service password-encryption
msk-ansusenko-gw-1(config)#username admin privilege 1 secret cisco
msk-ansusenko-gw-1(config)#ip domain-name ansusenko.rudn.edu
msk-ansusenko-gw-1(config)#crypto key generate rsa
msk-ansusenko-gw-1(config)#line vty 0 4
msk-ansusenko-gw-1(config-line)#transport input ssh
\end{minted}

\begin{center}
    \centering
    \includegraphics[width=\textwidth]{../images/02.png}
    \captionof{figure}{Первоначальная настройка маршрутизатора.}
    \label{02}
\end{center}

\subsection{Настройка Trunk-порта на коммутаторе}

На коммутаторе msk-ansusenko-sw-01 настроили порт 24 как trunk-порт для передачи трафика всех VLAN.

\begin{minted}{cisco}
msk-ansusenko-sw-01>enable
Password:
msk-ansusenko-sw-01#configure terminal
msk-ansusenko-sw-01(config)#interface fastEthernet 0/24
msk-ansusenko-sw-01(config-if)#switchport mode trunk
msk-ansusenko-sw-01(config-if)#switchport trunk allowed vlan all
msk-ansusenko-sw-01(config-if)#end
\end{minted}

После выполнения команды \texttt{switchport mode trunk} протокол линии на интерфейсе FastEthernet0/24 временно перешел в состояние down, а затем в up, что подтверждает успешную настройку trunk-порта.

\subsection{Настройка виртуальных интерфейсов на маршрутизаторе}

На маршрутизаторе msk-ansusenko-gw-1 настроили виртуальные интерфейсы, соответствующие номерам VLAN. Согласно таблице IP-адресов задали соответствующие IP-адреса (Рис. \ref{03}):

\begin{minted}{cisco}
msk-ansusenko-gw-1>enable
msk-ansusenko-gw-1#configure terminal
msk-ansusenko-gw-1(config)#interface fastEthernet 0/0
msk-ansusenko-gw-1(config-if)#no shutdown
msk-ansusenko-gw-1(config)#interface fastEthernet 0/0.2
msk-ansusenko-gw-1(config-subif)#encapsulation dot1Q 2
msk-ansusenko-gw-1(config-subif)#ip address 10.128.1.1 255.255.255.0
msk-ansusenko-gw-1(config-subif)#description management
msk-ansusenko-gw-1(config)#interface fastEthernet 0/0.3
msk-ansusenko-gw-1(config-subif)#encapsulation dot1Q 3
msk-ansusenko-gw-1(config-subif)#ip address 10.128.0.1 255.255.255.0
msk-ansusenko-gw-1(config-subif)#description servers
msk-ansusenko-gw-1(config)#interface fastEthernet 0/0.101
msk-ansusenko-gw-1(config-subif)#encapsulation dot1Q 101
msk-ansusenko-gw-1(config-subif)#ip address 10.128.3.1 255.255.255.0
msk-ansusenko-gw-1(config-subif)#description dk
msk-ansusenko-gw-1(config)#interface fastEthernet 0/0.102
msk-ansusenko-gw-1(config-subif)#encapsulation dot1Q 102
msk-ansusenko-gw-1(config-subif)#ip address 10.128.4.1 255.255.255.0
msk-ansusenko-gw-1(config-subif)#description departments
msk-ansusenko-gw-1(config)#interface fastEthernet 0/0.103
msk-ansusenko-gw-1(config-subif)#encapsulation dot1Q 103
msk-ansusenko-gw-1(config-subif)#ip address 10.128.5.1 255.255.255.0
msk-ansusenko-gw-1(config-subif)#description adm
msk-ansusenko-gw-1(config)#interface fastEthernet 0/0.104
msk-ansusenko-gw-1(config-subif)#encapsulation dot1Q 104
msk-ansusenko-gw-1(config-subif)#ip address 10.128.6.1 255.255.255.0
msk-ansusenko-gw-1(config-subif)#description other
\end{minted}

\begin{center}
    \centering
    \includegraphics[width=\textwidth]{../images/03.png}
    \captionof{figure}{Настройка виртуальных интерфейсов, соответствующих номерам VLAN.}
    \label{03}
\end{center}

Так же настроили trunk-порты. (Рис. \ref{05}):

\begin{center}
    \centering
    \includegraphics[width=\textwidth]{../images/05.png}
    \captionof{figure}{Настройка trunk-портов.}
    \label{05}
\end{center}

Для проверки выполнили команды (Рис. \ref{04}):

\begin{minted}{cisco}
msk-ansusenko-gw-1>show ip interface brief
msk-ansusenko-gw-1>show vlan
msk-ansusenko-gw-1>show interfaces trunk
\end{minted}

\textbf{Результат проверки:}
\begin{itemize}
    \item Физический интерфейс FastEthernet0/0: up/up
    \item Все суб-интерфейсы (VLAN 2,3,101,102,103,104) в состоянии up/up с правильными IP-адресами
    \item Команда \texttt{show vlan} на маршрутизаторе показывает только стандартные VLAN, что является нормой
    \item Команда \texttt{show interfaces trunk} не показывает информацию на маршрутизаторе
\end{itemize}

\begin{center}
    \centering
    \includegraphics[width=\textwidth]{../images/04.png}
    \captionof{figure}{Проверка заданных настроек.}
    \label{04}
\end{center}

\subsection{Проверка доступности устройств из разных VLAN}

\subsubsection{Проверка связи ПК dk-ansusenko-01 (VLAN 101)}

На ПК dk-ansusenko-01 (IP 10.128.3.2) выполнили проверку связи (Рис. \ref{06}):

\begin{enumerate}
    \item Пинг шлюза VLAN 101 (10.128.3.1):
    \begin{minted}{cisco}
    ping 10.128.3.1
    \end{minted}
    
    \item Пинг ПК из VLAN 102 (10.128.4.2):
    \begin{minted}{cisco}
    ping 10.128.4.2
    \end{minted}
    
    \item Пинг сервера (10.128.0.2):
    \begin{minted}{cisco}
    ping 10.128.0.2
    \end{minted}
\end{enumerate}

\begin{center}
    \centering
    \includegraphics[width=\textwidth]{../images/06.png}
    \captionof{figure}{Проверка связи с ПК dk-ansusenko-01.}
    \label{06}
\end{center}

\textbf{Результаты:}
\begin{itemize}
    \item Пинг шлюза: успешно (TTL=255) - связь в пределах VLAN работает
    \item Пинг ПК из VLAN 102: успешно (TTL=127) - меж VLAN маршрутизация работает
    \item Пинг сервера: успешно (TTL=127) - доступ к серверной ферме работает
\end{itemize}

\subsubsection{Проверка связи ПК из VLAN 102}

На ПК из VLAN 102 (IP 10.128.4.2) выполнили пинг своего шлюза (Рис. \ref{07}):

\begin{minted}{cisco}
ping 10.128.4.1
\end{minted}

\begin{center}
    \centering
    \includegraphics[width=\textwidth]{../images/07.png}
    \captionof{figure}{Проверка связи ПК из VLAN 102.}
    \label{07}
\end{center}

\subsubsection{Проверка связи между серверами}

На сервере ansusenko-web (10.128.0.2) выполнили пинг до сервера ansusenko-file (10.128.0.3) (Рис. \ref{08}):

\begin{minted}{cisco}
ping 10.128.0.3
\end{minted}

\begin{center}
    \centering
    \includegraphics[width=\textwidth]{../images/08.png}
    \captionof{figure}{Проверка связи между серверами.}
    \label{08}
\end{center}

\subsection{Изучение процесса передачи пакетов в режиме симуляции}

Для изучения процесса передачи ICMP-пакетов между разными VLAN использовали режим симуляции в Packet Tracer.

\subsubsection{Настройка симуляции}
\begin{enumerate}
    \item Переключились в режим \texttt{Simulation}
    \item Настроили фильтры только на ICMP и ARP (сняли все остальные)
    \item Запустили ping с ПК dk-ansusenko-01 (10.128.3.2) на ПК из VLAN 102 (10.128.4.2)
    \item Просмотрели наши пакеты после завершения симуляции (Рис. \ref{09}):
    
\begin{center}
    \centering
    \includegraphics[width=\textwidth]{../images/09.png}
    \captionof{figure}{Проверка связи между серверами.}
    \label{09}
\end{center}    
    
\end{enumerate}

\subsubsection{Наблюдения за движением пакета}

\textbf{Шаг 1: Пакет покидает ПК-источник (10.128.3.2)}
\begin{itemize}
    \item Ethernet-кадр без тега VLAN
    \item MAC-адрес назначения: MAC-адрес шлюза (полученный через ARP)
    \item IP-пакет: Source IP = 10.128.3.2, Destination IP = 10.128.4.2
\end{itemize}

\textbf{Шаг 2: Пакет поступает на коммутатор через access-порт}
\begin{itemize}
    \item Коммутатор добавляет тег VLAN 101 (802.1Q)
    \item Пакет направляется к trunk-порту (Fa0/24)
\end{itemize}

\textbf{Шаг 3: Пакет передается по trunk-порту}
\begin{itemize}
    \item В заголовке 802.1Q указан VLAN ID = 101
    \item Пакет поступает на маршрутизатор
\end{itemize}

\textbf{Шаг 4: Маршрутизация на маршрутизаторе}
\begin{itemize}
    \item Интерфейс Fa0/0.101 принимает пакет с тегом 101
    \item Маршрутизатор снимает тег, анализирует IP-адрес назначения
    \item Принимает решение о маршрутизации в сеть 10.128.4.0/24
    \item Добавляет новый тег VLAN 102
    \item Отправляет пакет через интерфейс Fa0/0.102
\end{itemize}

\textbf{Шаг 5: Пакет возвращается на коммутатор}
\begin{itemize}
    \item Коммутатор получает пакет с тегом 102 на trunk-порту
    \item Определяет порт назначения в VLAN 102
    \item Удаляет тег VLAN перед отправкой на access-порт
\end{itemize}

\textbf{Шаг 6: Пакет достигает ПК-назначения (10.128.4.2)}
\begin{itemize}
    \item ПК получает обычный Ethernet-кадр
    \item Формирует ICMP Echo Reply и отправляет обратно
\end{itemize}

\subsubsection{Содержимое пакетов (анализ PDU)}

\textbf{ICMP Echo Request (на выходе из источника):}
\begin{itemize}
    \item \textbf{Layer 2:} Source MAC = MAC ПК, Dest MAC = MAC шлюза
    \item \textbf{Layer 3:} Source IP = 10.128.3.2, Dest IP = 10.128.4.2, TTL = 128
    \item \textbf{Layer 4:} ICMP Type = 8 (Echo Request)
\end{itemize}

\textbf{ICMP Echo Request (на trunk-порту):}
\begin{itemize}
    \item \textbf{Layer 2:} Добавлен заголовок 802.1Q с VLAN ID = 101
    \item Остальные поля без изменений
\end{itemize}

\textbf{ICMP Echo Request (после маршрутизации):}
\begin{itemize}
    \item \textbf{Layer 2:} Source MAC = MAC маршрутизатора, Dest MAC = MAC ПК назначения
    \item \textbf{Layer 3:} TTL уменьшен до 127
    \item \textbf{Layer 2 заголовок:} 802.1Q с VLAN ID = 102
\end{itemize}

\subsection{Таблица распределения IP-адресов}

В соответствии с заданием, в сети использовались следующие IP-адреса (Таблица \ref{tab:ip}):

\begin{table}[h]
    \centering
    \caption{Таблица распределения IP-адресов}
    \label{tab:ip}
    \begin{tabular}{|p{4cm}|c|c|c|}
        \hline
        \textbf{Устройство} & \textbf{VLAN} & \textbf{IP-адрес} & \textbf{Назначение} \\
        \hline
        msk-ansusenko-gw-1 & - & 10.128.0.1 & Шлюз серверной фермы \\
        \hline
        msk-ansusenko-gw-1 & - & 10.128.1.1 & Шлюз управления \\
        \hline
        msk-ansusenko-gw-1 & - & 10.128.3.1 & Шлюз ДК \\
        \hline
        msk-ansusenko-gw-1 & - & 10.128.4.1 & Шлюз кафедр \\
        \hline
        msk-ansusenko-gw-1 & - & 10.128.5.1 & Шлюз администрации \\
        \hline
        msk-ansusenko-gw-1 & - & 10.128.6.1 & Шлюз других пользователей \\
        \hline
        ansusenko-web & 3 & 10.128.0.2 & Web-сервер \\
        \hline
        ansusenko-file & 3 & 10.128.0.3 & Файловый сервер \\
        \hline
        ansusenko-mail & 3 & 10.128.0.4 & Почтовый сервер \\
        \hline
        dk-ansusenko-01 & 101 & 10.128.3.2 & ПК пользователя (ДК) \\
        \hline
        dep-ansusenko-1 & 102 & 10.128.4.2 & ПК пользователя (кафедра) \\
        \hline
        adm-ansusenko-1 & 103 & 10.128.5.2 & ПК пользователя (администрация) \\
        \hline
        other-ansusenko-1 & 104 & 10.128.6.2 & ПК пользователя (другие) \\
        \hline
    \end{tabular}
\end{table}

\section{Выводы}

В результате выполнения лабораторной работы мы построили статическую маршрутизацию VLAN и убедились, что она настроена корректно. Все устройства в разных VLAN имеют возможность взаимодействовать друг с другом через маршрутизатор, при этом трафик изолирован на втором уровне модели OSI и маршрутизируется на третьем уровне.

\end{document}
